\documentclass[../../../include/open-logic-section]{subfiles}

\begin{document}

\olfileid[pt]{sth}{choice}{justifications}
\olsection{Considerações Intrínsecas sobre a Escolha}

A questão mais ampla, então, é se a Boa Ordenação, ou a Escolha, ou
de fato a comparabilidade de todos os conjuntos quanto ao tamanho---não
importa qual---podem ser justificadas. 

Eis uma tentativa de justificação \emph{intrínseca}. De volta a
\olref[z][story]{sec}, introduzimos vários princípios
sobre a hierarquia. Um deles merece ser repetido:
\begin{enumerate}
	\item[] \stagesacc. Para qualquer estágio $S$, e para quaisquer conjuntos que foram
	formados \emph{antes} do estágio $S$: forma-se no estágio $S$ um conjunto cujos
	membros são exatamente esses conjuntos. Nada mais se forma no
	estágio~$S$. 
\end{enumerate}
De fato, muitos autores sugeriram que o Axioma da Escolha pode ser
justificado via (algo como) este princípio. Daremos brevemente um
resumo dessa abordagem.

Começamos por um resultado simples, que oferece \emph{mais uma}
formulação equivalente da Escolha:

\begin{thm}[em $\ZF$]\ollabel{choiceset}
A Escolha é equivalente ao princípio seguinte. Se os !!{element}s
de $A$ são disjuntos dois a dois e não vazios, então existe algum $C$ tal que $C
\cap x$ é um singular para todo o $x \in A$. (Chamamos a tal $C$ um
{conjunto de escolha} para $A$.)
\end{thm}

A prova deste resultado é direta, e a deixamos como
exercício para o leitor. 

\begin{prob}
Prove \olref[sth][choice][justifications]{choiceset}. Se tiver dificuldades,
encontrará uma prova em \cite[pp.~242--3]{Potter2004}.
\end{prob}

O ponto essencial é que um conjunto de escolha para $A$ é precisamente o contradomínio de
uma função de escolha para $A$. Assim, para justificar a Escolha, podemos simplesmente tentar
justificar a sua formulação equivalente, em termos da existência de
conjuntos de escolha. E é exatamente isso que faremos agora. 

Sejam os !!{element}s de $A$ disjuntos dois a dois e não vazios. Por \stageshier{}
(ver \olref[z][story]{sec}), $A$ é formado em algum estágio~$S$. Note-se
que todos os !!{element}s de $\bigcup A$ estão disponíveis antes do estágio
$S$. Agora, por \stagesacc{}, para \emph{quaisquer} conjuntos formados
antes de~$S$, forma-se um conjunto cujos membros são exatamente esses conjuntos.
Dito de outro modo: toda a \emph{possível} coleção de conjuntos disponíveis mais cedo
existirá em~$S$. Mas é certamente \emph{possível} seleccionar
objetos que possam ser reunidos num conjunto de escolha para~$A$; isso é apenas
um certo subconjunto muito específico de $\bigcup A$. Logo: existe tal conjunto de escolha,
como requerido.

Bem, esta é uma tentativa \emph{muito} rápida de oferecer uma justificação da
Escolha por razões intrínsecas. Mas, para aprofundar esta ideia, deve ler
o desenvolvimento elegante de Potter (\citeyear[\S14.8]{Potter2004}).

\end{document}
